
\clearpage
\onecolumn
\section*{Appendix: Use of AI/LLM Tools}
\addcontentsline{toc}{section}{Appendix: Use of AI/LLM Tools}

ChatGPT was used as a limited supporting tool during this project. The tool was used during parts of the project work in 2026 through the ChatGPT web interface. The implementation, numerical experiments, figures, tables, and final scientific interpretation are my own responsibility.

For the code, ChatGPT was used mainly for support while developing and cleaning the project. This included help with discussing repository organization, interpreting terminal errors, checking Git and \LaTeX{} commands, improving file names and structure, and making comments and code organization easier to read. I also used ChatGPT for debugging support and for suggestions about useful consistency checks. All code suggestions were reviewed before use, and the submitted code was tested by running the project workflow and unit tests.

For the written report, ChatGPT was used as a language and feedback tool. It helped with proofreading, improving wording, checking notation consistency, and giving feedback on whether the report addressed the course evaluation criteria. I reviewed and edited the suggestions before including them in the final text.

\noindent\textbf{Table A1 gives the level of AI assistance used in the written report.}

\begin{center}
\begin{tabular}{|p{0.25\textwidth}|p{0.10\textwidth}|p{0.55\textwidth}|}
\hline
\textbf{Report section} & \textbf{Level} & \textbf{Description} \\
\hline
Abstract & 1 & Language and clarity check. \\
Introduction & 1 & Minor wording and structure feedback. \\
Theory and methods & 1 & Language cleanup and notation consistency. \\
Implementation and validation & 1 & Wording and readability feedback. \\
Results and discussion & 1 & Language cleanup and feedback on clarity of the discussion. \\
Conclusion & 1 & Minor wording improvements. \\
AI declaration & 2 & Edited to follow the course LLM declaration guidelines. \\
\hline
\end{tabular}
\end{center}

For text assistance, level 0 means no LLM assistance, level 1 means language-only assistance, level 2 means editorial assistance, and level 3 means generative drafting. In this report the assistance was mainly level 1, with some level 2 editorial feedback for the AI declaration and general report structure. No report section was included without my own review and approval.

\noindent\textbf{Table A2 gives the level of AI assistance used for code and workflow.}

\begin{center}
\begin{tabular}{|p{0.25\textwidth}|p{0.10\textwidth}|p{0.55\textwidth}|}
\hline
\textbf{Component} & \textbf{Level} & \textbf{Description} \\
\hline
Repository organization & 1 & Suggested cleaner folder structure and Git commands. \\
Terminal and Git workflow & 1 & Help interpreting terminal errors and checking commands. \\
Plotting scripts & 1 & Help with readability, figure sizing, and report-oriented plotting. \\
Python source files & 1--2 & Debugging support, comments, naming, and short code suggestions. \\
Tests & 1--2 & Suggestions for useful consistency checks. \\
\LaTeX{}/report formatting & 1 & Help with figure placement, formatting, and compilation commands. \\
\hline
\end{tabular}
\end{center}

For code assistance, level 0 means no LLM assistance, level 1 means debugging or explanation assistance, level 2 means short snippets or small blocks, level 3 means skeleton generation, and level 4 means substantial file-level generation. In this project the assistance was mainly level 1--2: debugging support, command checking, code cleanup suggestions, comments, naming, and short code suggestions. ChatGPT was not used to fabricate numerical results, choose final data, or replace my own calculations and scientific interpretation.

All reported values were generated by running the project code in the GitHub repository. The numerical results were checked against exact diagonalization, analytic limits, unit tests, or other relevant consistency checks. I take responsibility for the submitted code, figures, tables, numerical values, and scientific interpretation, and I can explain the equations, algorithms, and results used in the final submission.
